% TOP_PROBLEM: {"id": 2306117, "title": "Research Problems in Function Theory — Problem 6.117", "category_id": 8, "category": {"id": 8, "name": "analysis", "display_name": "Analysis", "description": "Limits, continuity, calculus, and function theory.", "slug": "analysis", "order_index": 8, "created_at": "2026-07-31T15:26:25.671Z"}, "status": "open"}
% FIRST_POSTED: null
% ENTRY_KIND: statement_only
% Imported from: batch04.tex; UnsolvedMath ID 2306117
\documentclass[11pt]{article}
\usepackage[margin=25mm]{geometry}
\usepackage{fontspec}
\setmainfont{Latin Modern Roman}
\usepackage{amsmath,amssymb,mathtools}
\usepackage{unicode-math}
\setmathfont{Latin Modern Math}
\usepackage{xurl,hyperref}
\hypersetup{colorlinks=true,urlcolor=blue}
\setcounter{secnumdepth}{0}
\setlength{\parindent}{0pt}
\setlength{\parskip}{5pt}
\setlength{\emergencystretch}{3em}
\newcommand{\sref}[2]{\hyperlink{source:#1:#2}{[S#2]}}
\newcommand{\cataloguescope}[1]{\paragraph{Recorded target.}#1}
\begin{document}
% UnsolvedMath ID 2306117; source catalogue code AMR-022-6117
\setcounter{section}{2306116}
\section{Research Problems in Function Theory — Problem 6.117}\label{2306117}
{\small\sffamily UnsolvedMath ID 2306117 · Analysis}
\subsection{Definitions and mathematical statement}

\paragraph{Shared notation.}$\mathbb N=\{1,2,\ldots\}$, and $\mathbb Z,\mathbb Q,\mathbb R,\mathbb C$ denote the integers, rationals, reals and complex numbers. Any other conventions are specified locally.


For a proper plane domain $U$, a Borel set $E\subseteq\partial U$, and $z\in U$, $\omega(z,E,U)$ denotes harmonic measure: the probability that planar Brownian motion begun at $z$ first exits $U$ through $E$. For an arbitrary set, use the associated outer measure; for sets off the boundary, use their intersection with $\partial U$.
For an open connected $G\subseteq\mathbb C$ containing zero, axial symmetry means $[z,\overline z]\subseteq G$ for every $z\in G$. For $t>0$, let $G_t^h$ and $G_t^v$ be the zero-containing components of $G\cap\{\operatorname{Im}z<t\}$ and $G\cap\{\operatorname{Re}z<t\}$. Define
\[h_G(t)=\omega(0,\partial G_t^h\cap\{\operatorname{Im}z=t\},G_t^h),\qquad
v_G(t)=\omega(0,\partial G_t^v\cap\{\operatorname{Re}z=t\},G_t^v).\]
\paragraph{Statement recorded in the supplied source.}
For axially symmetric domains $D_1,D_2$ containing zero, does each of the following separately imply $D_1=D_2$: $h_{D_1}(t)=h_{D_2}(t)$ for every $t>0$; $v_{D_1}(t)=v_{D_2}(t)$ for every $t>0$; or both equalities together? If either individual implication holds, how few values of $t$ suffice: for example, is an infinite set of values whose corresponding horizontal or vertical lines meet both domains enough, or is a dense set required? For each implication that holds, does it remain valid for all simply connected domains containing zero, or for all domains containing zero without that restriction? \sref{2306117}{2}
\subsection{Short English statement}
Can horizontal or vertical harmonic-measure profiles identify a domain, how many samples suffice, and can the symmetry assumption be weakened or removed?
\subsection{Sources}
\begin{description}
\item[\hypertarget{source:2306117:1}{S1}] ULAM.AI and the UnsolvedMath Contributors, UnsolvedMath: A Curated Collection of Open Mathematics Problems, record 2306117 (AMR-022-6117). CC BY 4.0. \url{https://huggingface.co/datasets/ulamai/UnsolvedMath}.
\item[\hypertarget{source:2306117:2}{S2}] W. K. Hayman and E. F. Lingham, Research Problems in Function Theory (2018), Problem 6.117, complete problem and update. \url{https://arxiv.org/abs/1809.07200}.
\end{description}
\label{end:2306117}
\end{document}
